\documentclass[11pt,a4paper]{article}
\usepackage{../pelm}
\begin{document}
\paper{Quiz 2}{Weeks 3--4: How LLMs Work, and Principles of Effective Prompting}{10}{10 minutes}

\begin{questions}

\q{A \textbf{token} is best described as:}
  {Always a single word}
  {A piece of text --- a word, part of a word, a character, or punctuation}
  {A single letter}
  {A complete sentence}

\q{Why do modern systems break words into sub-word pieces rather than using a fixed list of whole
   words?}
  {It makes the model smaller}
  {It allows any word, including unseen ones, to be assembled from familiar pieces}
  {It improves the model's grammar}
  {It reduces the training time}

\q{Why does the same content typically cost more tokens in Turkish than in English?}
  {Turkish sentences are longer}
  {Turkish is agglutinative, and stacked suffixes fragment into several pieces}
  {Turkish uses a different alphabet}
  {Turkish text is processed twice}

\q{An assistant miscounts how many times a letter appears in a long word. The best explanation is:}
  {The model is not good at mathematics}
  {The model never perceived individual letters --- it works on tokens}
  {The context window was exceeded}
  {The temperature setting was too high}

\q{Which statement best describes the model's core operation?}
  {It retrieves the most relevant stored fact}
  {It produces a probability for every possible next token, then one is chosen}
  {It searches the internet and summarises the results}
  {It applies a set of rules written by its developers}

\q{Why does the same prompt sometimes produce different answers?}
  {The model is retrained between requests}
  {There is deliberate randomness in choosing which token comes next}
  {The context window changes size}
  {The model remembers the previous answer and avoids repeating it}

\q{Asking the same question twice and receiving a similar answer tells you:}
  {The answer is correct}
  {The pattern is strongly represented in the training data}
  {The model is confident}
  {The context window is large enough}

\q{Which statement about hallucination is correct?}
  {It occurs when the context window overflows}
  {It is a bug that will be removed in future versions}
  {There is no step at which the system checks whether what it says is true}
  {It only happens with small models}

\q{A colleague adds ``Be accurate and only use reliable sources'' to a prompt about a topic the
   model knows little about. What is the most likely effect?}
  {The answer becomes accurate}
  {The model refuses to answer}
  {The answer sounds more authoritative without becoming better informed}
  {The context window expands}

\q{Output arrives with correct content but as a wall of unstructured text when you wanted three
   short bullets. According to the Week 4 diagnosis, which lever should you pull?}
  {Context}
  {Task}
  {Format}
  {Examples}

\end{questions}
\end{document}
